\documentclass{article}
\usepackage{amsmath}
\usepackage{graphicx}
\usepackage{geometry}
\geometry{a4paper}

\title{Diseño del Dispositivo IoT para el Pesaje de Palets de Queso y Cajas de Cartón}
\author{Ingeniería Electrónica}
\date{\today}

\begin{document}

\maketitle

\section{Introducción}
Este documento presenta el diseño de un sistema IoT para el pesaje en tiempo real de palets de queso y cajas de cartón en una empresa láctea. El objetivo es anticiparse al vaciado de los palets mediante el uso de celdas de carga y sensores IoT, lo que optimiza la gestión de recursos y reduce los tiempos de inactividad en la línea de producción.

\section{Diagrama de Bloques del Sistema}
El sistema se compone de los siguientes bloques principales: 

\begin{center}
\includegraphics[width=\textwidth]{diagrama_de_bloques.png} % Asegúrate de incluir una imagen del diagrama de bloques si la tienes
\end{center}

\section{Componentes del Diseño}

\subsection{A. Sensores de Peso (Celdas de Carga)}
Las celdas de carga son dispositivos que convierten la fuerza de compresión (el peso) en una señal eléctrica. Cada palet de queso o caja de cartón estará equipado con 4 celdas de carga.

\begin{itemize}
    \item \textbf{Selección del Sensor:} Se utilizarán celdas de carga tipo S (tensión), que son comúnmente usadas en aplicaciones industriales debido a su alta precisión y fiabilidad.
    \item \textbf{Interfaz:} Las celdas de carga estarán conectadas al amplificador de señal HX711, que amplifica la señal de salida para que pueda ser leída por el microcontrolador.
\end{itemize}

\subsection{B. Amplificador de Señal}
Las celdas de carga proporcionan una señal analógica de muy baja intensidad. Para que el microcontrolador pueda leer esta señal, se utiliza un amplificador de señal.

\begin{itemize}
    \item \textbf{Amplificador de Señal:} Se utilizará el amplificador de señal \textbf{HX711}, que convierte la señal analógica de las celdas de carga a una señal digital legible por el microcontrolador.
\end{itemize}

\subsection{C. Microcontrolador / Procesador}
El microcontrolador es responsable de leer las señales de peso proporcionadas por las celdas de carga y enviarlas a la plataforma de monitoreo.

\begin{itemize}
    \item \textbf{Microcontrolador recomendado:} Se utilizará un microcontrolador \textbf{ESP32} o \textbf{ESP8266}, que soportan conectividad WiFi integrada, ideal para el envío de datos a la plataforma en tiempo real.
    \item \textbf{Funciones del microcontrolador:} El microcontrolador leerá los datos del HX711, calculará el peso y enviará esta información a la plataforma de monitoreo.
\end{itemize}

\subsection{D. Módulo de Comunicación Inalámbrica}
El sistema IoT debe enviar los datos de peso a la plataforma de monitoreo en tiempo real. Para ello, se utilizarán módulos de comunicación inalámbrica.

\begin{itemize}
    \item \textbf{WiFi:} Si la infraestructura de la planta lo permite, se utilizará el módulo WiFi integrado en el ESP32 para enviar los datos a la red local.
    \item \textbf{Alternativa - LoRa:} En caso de que no haya cobertura WiFi suficiente, se puede utilizar LoRa para la transmisión de datos a larga distancia con bajo consumo de energía.
    \item \textbf{GSM:} Si se requiere cobertura en zonas alejadas de la planta, se puede integrar un módulo GSM para enviar los datos mediante red celular.
\end{itemize}

\subsection{E. Fuente de Alimentación}
El dispositivo IoT necesitará una fuente de alimentación confiable.

\begin{itemize}
    \item \textbf{Fuente de alimentación:} Se utilizará una batería recargable de tipo Li-Ion o Li-Po, complementada con un circuito de carga para garantizar un funcionamiento autónomo.
    \item \textbf{Alternativa de energía:} En caso de que sea posible, el sistema puede conectarse a una fuente de energía AC/DC para evitar depender de la batería.
\end{itemize}

\subsection{F. Plataforma de Monitoreo / ERP}
Los datos de peso se enviarán a una plataforma centralizada para su monitoreo.

\begin{itemize}
    \item \textbf{Plataforma de monitoreo:} Se desarrollará una interfaz visual, accesible desde una aplicación web o móvil, que permitirá a los supervisores monitorear en tiempo real el peso de los palets.
    \item \textbf{Integración con ERP:} La plataforma de monitoreo también estará integrada con el sistema ERP de la empresa, lo que permitirá gestionar los inventarios y los recursos logísticos de manera eficiente.
\end{itemize}

\section{Diseño de Hardware}

\subsection{Conexiones entre los Componentes}
A continuación, se detallan las conexiones necesarias para que los componentes interactúen de manera efectiva:

\begin{itemize}
    \item Las celdas de carga se conectan al amplificador HX711 a través de 4 cables: VCC, GND, DT (data) y SCK (clock).
    \item El amplificador HX711 se conecta al microcontrolador ESP32 mediante los pines DT (Data) y SCK (Clock).
    \item El microcontrolador ESP32 se conecta a la red WiFi o LoRa para enviar los datos a la plataforma de monitoreo.
\end{itemize}

\subsection{Esquema Básico de Conexiones}
El esquema básico de conexiones es el siguiente:

\[
\text{Celdas de Carga} \rightarrow \text{Amplificador HX711} \rightarrow \text{Microcontrolador ESP32} \rightarrow \text{Plataforma de Monitoreo}
\]

\section{Consideraciones de Software}

\subsection{Lectura del Peso}
El software del microcontrolador debe leer las señales analógicas de las celdas de carga y convertirlas a valores digitales mediante el uso del amplificador HX711. La lectura se realiza periódicamente para obtener el peso total de cada palet.

\subsection{Transmisión de Datos}
Una vez calculado el peso, el microcontrolador enviará los datos a la plataforma de monitoreo mediante WiFi o LoRa.

\subsection{Alertas}
El sistema debe generar alertas automáticas cuando el peso del palet se acerque al umbral crítico, notificando a los operarios para que procedan con la reposición del material.

\section{Pasos para la Implementación}

\begin{enumerate}
    \item Desarrollar el software para leer las señales de las celdas de carga y calcular el peso de los palets.
    \item Implementar la comunicación del microcontrolador con la plataforma de monitoreo.
    \item Realizar pruebas de calibración con diferentes pesos de palets.
    \item Integrar el sistema con el ERP de la empresa para la gestión de inventarios.
    \item Implementar el sistema en toda la planta y evaluar su rendimiento.
\end{enumerate}

\section{Conclusión}
El sistema IoT propuesto para el pesaje de palets de queso y cajas de cartón permitirá optimizar el flujo de materiales y reducir los tiempos muertos durante el proceso de embalaje. Con la implementación de un sistema de pesaje y monitoreo en tiempo real, se mejorará la eficiencia operativa y la planificación logística.

\end{document}

